\section*{Bab \ref{ch:b3:developmental-genetics} --- Genetika Perkembangan dan Evolusi Bentuk}

\begin{solution}{exo:b3:developmental-genetics:1}
Pengaruh induk: \emph{bicoid}, dengan embrio dari induk mutan tidak punya kepala
dan toraks. Senjang: \emph{Krüppel}, dengan sebuah blok segmen yang bersambungan
hilang. Aturan pasangan: \emph{fushi tarazu} (atau \emph{even-skipped}), dengan tiap
segmen berselang hilang. Kekutuban segmen: \emph{hedgehog} (atau
\emph{engrailed}), dengan sebagian tiap segmennya digantikan sebuah bayangan cermin
sisanya.
\end{solution}

\begin{solution}{exo:b3:developmental-genetics:2}
Sebuah gen yang produknya ditaruh di dalam telurnya oleh induknya, sehingga
fenotipe embrionya mengikuti genotipe induknya. Embrionya tidak punya kepala dan
toraks lalu mengusung struktur belakang pada kedua ujungnya; mereka mati.
Alel liar dari ayahnya ada di dalam zigotnya, tetapi mRNA \emph{bicoid}
dibuat selama oogenesisnya oleh sel perawat induknya lalu ditambatkan pada
kutubnya sebelum pembuahannya; salinan zigotnya sendiri akan ditranskripsi
berjam-jam kemudian, sesudah waktu landaiannya berlalu.
\end{solution}

\begin{solution}{exo:b3:developmental-genetics:3}
\omterm{def:b3:developmental-genetics:hox}{Kolinearitas}: urutan \omterm{def:b3:developmental-genetics:hox}{gen Hox} sepanjang kromosomnya memadani
urutan ranah pengungkapannya sepanjang tubuhnya. Keunggulan belakang:
tempat beberapa di antaranya diungkapkan, yang paling belakanglah yang menetapkan jati dirinya.
Seekor mencit tanpa \emph{Hoxc8} menunjukkan sebuah perubahan ke arah depan --- ruas
tulang belakang pinggang pertamanya mengambil jati diri sebuah ruas toraks lalu mengusung
sebuah rusuk. \emph{Antennapedia} di kepalanya mengubah sungutnya menjadi kaki, yakni
jati diri segmen toraks keduanya.
\end{solution}

\begin{solution}{exo:b3:developmental-genetics:4}
Sebuah bibir punggung yang dicangkokkan mengimbas sebuah sumbu kedua yang lengkap yang terbuat kebanyakan dari
sel inangnya: cangkoknya memerintah tetangganya --- yakni sebuah \omterm{def:b3:developmental-genetics:organiser}{penata}. Molekulnya
adalah penentang BMP yang disekresikan (Chordin, Noggin, Follistatin).
Ektodermanya menjadi saraf ketika pengisyaratan BMP-nya tidak ada dan menjadi epidermis
ketika BMP-nya hadir, sehingga saraf adalah apa yang dikerjakan ektodermanya ketika tidak diberi tahu apa-apa:
\omterm{def:b3:developmental-genetics:organiser}{penatanya} bekerja dengan membungkam sebuah perintah.
\end{solution}

\begin{solution}{exo:b3:developmental-genetics:5}
$k = \ln 2/(35\times 60) = \qty{3.3e-4}{s^{-1}}$; $\lambda = \sqrt{4/
3.3\times 10^{-4}} = \qty{110}{\micro m}$; nisbahnya $\mathrm{e}^{250/110}
= \mathrm{e}^{2.27} = 9.7$.
\end{solution}

\begin{solution}{exo:b3:developmental-genetics:6}
Empat salinan melipatduakan $c_{0}$: batasnya bergerak ke belakang sebesar $\lambda\ln 2 =
\qty{69}{\micro m}$ ke \qty{309}{\micro m}; satu salinan membagi duanya: ke depan
sebesar \qty{69}{\micro m} ke \qty{171}{\micro m}. Tiap geserannya
\qty{14}{\%} panjang telurnya --- lebih besar daripada yang teramati, yang dengan sendirinya menjadi
bukti bagi mekanisme yang meredam kebergantungan dosisnya.
\end{solution}

\begin{solution}{exo:b3:developmental-genetics:7}
$\delta x = \lambda\,\delta c/c$: bagi \qty{8}{\micro m}, $\delta c/c =
8/100 = \qty{8}{\%}$. Sebuah cacahan Poisson $N$ punya galat nisbi
$1/\sqrt{N}$, sehingga $N = 1/0.08^{2} \approx 160$ molekul harus dicacah
--- yakni beberapa persen molekul yang hadir di dalam sebuah inti, atau sebuah cacahan
tunggal yang diintegralkan sepanjang waktu.
\end{solution}

\begin{solution}{exo:b3:developmental-genetics:8}
$k = \ln 2/35 = \qty{0.0198}{min^{-1}}$: $1 - \mathrm{e}^{-kt}$ memberi
$0.70$ pada \qty{60}{min}, $0.83$ pada \qty{90}{min}, $0.91$ pada
\qty{120}{min}, $0.95$ pada \qty{150}{min}. Landaiannya masih naik
beberapa persen selagi dibaca, yang akan menggeser sebuah batas sebesar
$\lambda\ln(1/0.9) \approx \qty{10}{\micro m}$, yakni sekitar satu inti, sepanjang
jendela pembacaannya: embrionya harus entah membaca pada sebuah waktu yang tetap atau memakai
sebuah pembacaan yang tak peka pada aras keseluruhannya.
\end{solution}

\begin{solution}{exo:b3:developmental-genetics:9}
Sebuah ZPA kedua memberi sebuah penggandaan cermin, 4-3-2-2-3-4: sel
depannya kini melihat Shh yang tinggi lalu menjadi jari belakang. Sebuah manik
yang melepaskan sedikit Shh memberi sebuah penggandaan sebagian --- yakni sebuah jari 2 depan
tambahan, katakanlah 2-2-3-4 --- sebab sel depannya melihat hanya kepekatan
rendah yang menentukan jari 2.
\end{solution}

\begin{solution}{exo:b3:developmental-genetics:10}
Penyelesaian khususnya $s/k$; penyelesaian homogennya $\cosh(x/\lambda)$,
$\sinh(x/\lambda)$; tanpa fluks pada $0$ mematikan $\sinh$-nya; $c(L) = 0$ menetapkan
amplitudonya: $c(x) = (s/k)\bigl[1 - \cosh(x/\lambda)/\cosh(L/\lambda)\bigr]$.
Tampangnya sebuah dataran dekat $x = 0$ yang jatuh ke nol di dalam sebuah lapisan
batas selebar $\lambda$ pada sumurnya --- bukan sebuah eksponensial dari sebuah
titik, melainkan sebuah medan yang rata dengan sebuah tepi; kepekatan tertingginya paling jauh
dari sumurnya. Bagi $L \ll \lambda$, $\cosh u \approx 1 +
u^{2}/2$ dan $c \approx s(L^{2} - x^{2})/(2D)$: yakni sebuah parabola yang tak bergantung
pada $k$, sebab molekulnya mencapai sumurnya sebelum dapat meluruh.
\end{solution}

\begin{solution}{exo:b3:developmental-genetics:11}
$x_{1} = \lambda\ln(1/0.3) = 1.20\lambda$ dan $x_{2} = \lambda\ln(1/0.08)
= 2.53\lambda$: daerah pertamanya selebar $1.20\lambda$, yang kedua
$1.32\lambda$, keduanya tak bergantung pada $c_{0}$, yang menggesernya bersama-sama.
Daerah ketiganya, $L - 2.53\lambda$, menyerap semua keragaman panjang
telurnya: pada sebuah telur \qty{450}{\micro m} dengan $\lambda = \qty{100}{\micro m}$
abdomennya \qty{197}{\micro m}, pada sebuah telur \qty{550}{\micro m}
\qty{297}{\micro m} --- polanya tidak terskala. Mekanismenya: sebuah landaian
kedua dari kutub belakangnya (Nanos, Caudal, atau sistem ujungnya) yang dibaca
bersama Bicoidnya, sehingga kedudukannya dipancangkan nisbahnya; atau sebuah molekul
``pemuai'' yang dibuat tempat morfogennya rendah yang memanjangkan $\lambda$ sampai landaiannya mengisi telurnya.
\end{solution}

\begin{solution}{exo:b3:developmental-genetics:12}
Sebuah perubahan penyandi mengubah proteinnya di mana pun ia bekerja: \emph{Pitx1} juga
membangun \omterm{def:b3:endocrinology:axes}{hipofisis} dan rahangnya, \emph{\omterm{def:b3:developmental-genetics:organiser}{sonic hedgehog}} juga memola
tabung saraf, usus, dan giginya, sehingga sebuah protein yang rusak mematikan atau
melumpuhkan secara pleiotropis. Sebuah penggalak menggerakkan pengungkapan di satu jaringan
pada satu waktu; menghapusnya menyingkirkan struktur yang akan dibangun jaringan itu
lalu membiarkan tiap kegunaan lain proteinnya utuh. Penggalaknya
modular dan kerap sebagian berlebihan, sehingga langkah antaranya dapat hidup,
dan jalan yang sama --- sebuah sakelar yang hilang --- telah ditempuh secara mandiri di
banyak danau ikan duri.
\end{solution}

\begin{solution}{pb:b3:developmental-genetics:1}
\textbf{1.} $k = \ln 2/(35\times 60) = \qty{3.3e-4}{s^{-1}}$; $1/k =
\qty{3030}{s} = \qty{50}{min}$.
\textbf{2.} $\lambda = \sqrt{4/3.3\times 10^{-4}} = \qty{110}{\micro m}$;
telurnya $4.5\lambda$.
\textbf{3.} $\mathrm{e}^{4.5} \approx 90$ dari kutub ke kutub;
$\mathrm{e}^{8.5/110} = 1.08$, yakni sebuah perubahan \qty{8}{\%} tiap inti.
\textbf{4.} $J = 50\times\sqrt{4\times 3.3\times 10^{-4}} = 50\times
0.036 = \qty{1.8}{nM\,\micro m\,s^{-1}}$. Satu nM adalah $0.6$ molekul
tiap \unit{\micro m^3}, sehingga sekitar $1.1$ molekul tiap \unit{\micro m^2}
tiap detik.
\textbf{5.} Volume intinya $\tfrac{4}{3}\pi\times 27 = \qty{113}{\micro
m^3}$. Pada \qty{50}{nM}, $30$ molekul tiap \unit{\micro m^3}:
yakni \num{3400} molekul. Di tengah telurnya $c = 50\,\mathrm{e}^{-2.27} =
\qty{5.2}{nM}$: yakni sekitar \num{350} molekul.
\textbf{6.} $kt = 1.19$ pada \qty{60}{min}: $0.70$; $kt = 2.38$ pada
\qty{120}{min}: $0.91$. Landaiannya dalam \qty{10}{\%} dari keadaan mantapnya
ketika dibaca --- masih naik sebesar jumlah yang akan memindahkan sebuah batas sejauh
$\lambda\ln(1/0.91) \approx \qty{10}{\micro m}$, yakni sebuah inti.
\textbf{7.} $x_{1} = 110\ln(1/0.3) = \qty{132}{\micro m}$
(\qty{26}{\%}); $x_{2} = 110\ln(1/0.08) = \qty{278}{\micro m}$
(\qty{56}{\%}).
\textbf{8.} Empat salinan: $c_{0} = \qty{100}{nM}$; kedua batasnya bergerak
ke belakang sebesar $\lambda\ln 2 = \qty{76}{\micro m}$, ke $208$ dan
\qty{354}{\micro m}. Enam salinan: $\lambda\ln 3 = \qty{121}{\micro m}$,
ke $253$ dan \qty{399}{\micro m}.
\textbf{9.} Satu salinan: ke depan sebesar \qty{76}{\micro m}, ke $56$ dan
\qty{202}{\micro m}. Daerah kepalanya menyusut dari $132$ menjadi
\qty{56}{\micro m}; toraksnya menjaga lebarnya (\qty{146}{\micro m}) lalu
bergerak ke depan; abdomennya tumbuh dari $222$ menjadi \qty{298}{\micro m}.
\textbf{10.} $\delta x = 110\times 0.1 = \qty{11}{\micro m}$, yakni $1.3$
jarak inti.
\textbf{11.} Satu inti: $\delta c/c = 8.5/110 = \qty{7.7}{\%}$. Dari
\qty{10}{\%}, $n = (10/7.7)^{2} \approx 1.7$: yakni dua pembacaan bebas
(atau rata-rata dua inti yang bertetangga) sudah mencukupi.
\textbf{12.} $1/\sqrt{350} = \qty{5.3}{\%}$: sebuah cacahan seketika yang tunggal
sudah lebih cermat daripada \qty{10}{\%} yang dianggap, sehingga derau
yang teramati bukanlah derau pencacahan saja melainkan datang dari kinetika
pengikatan dan mesin pembacaannya --- dan merata-rata sepanjang waktu
memperbaikinya.
\textbf{13.} $278/450 = \qty{62}{\%}$ dan $278/550 = \qty{51}{\%}$ panjang
telurnya: yakni sebuah selisih \qty{11}{\%}, enam inti. Polanya tidak
terskala dengan $\lambda$ yang terpancang.
\textbf{14.} Bila $D \propto L^{2}$ dengan $k$ terpancang, $\lambda \propto L$ dan
$x_{T}/L = (\lambda/L)\ln(c_{0}/T)$ sama pada tiap telurnya: yakni penskalaan yang
sempurna. Tetapi $D$ adalah sifat molekulnya dan sitoplasmanya, bukan sifat
ukuran telurnya, sehingga ini tidak masuk akal sebagaimana dinyatakan; penskalaan pada embrio
yang nyata bersifat sebagian dan datang dari pembacaan lain.
\textbf{15.} $c^{5}/(c^{5} + T^{5}) = 0.1$ pada $c/T = 9^{-1/5} = 0.64$ dan
$0.9$ pada $c/T = 9^{1/5} = 1.55$: yakni sebuah faktor $2.4$ pada kepekatannya,
$\Delta x = 110\ln 2.4 = \qty{97}{\micro m}$ --- yakni sebelas inti, jauh
lebih lebar daripada batas yang teramati yang dua atau tiga inti; pembacaan landaiannya
saja tidak cukup tajam.
\textbf{16.} Kecuramannya memetakan sebuah kepekatan tertentu menjadi menyala atau mati dengan lebih
tegas, sehingga daerah peralihannya lebih sempit; tetapi sebuah galat \qty{10}{\%}
pada kepekatannya tetap memindahkan titik tempat ambangnya
dilintasi sejauh $\lambda\,\delta c/c$: tanggapannya dapat berupa sebuah tangga yang sempurna
dan tangganya tetap mendarat di tempat yang salah.
\textbf{17.} Bila keduanya menyala, masing-masing akan menekan yang lain, sehingga
sedikitnya satu akan terdorong mati; dengan penekanan yang kooperatif, satu-satunya
keadaan mantapnya adalah satu-menyala-satu-mati, dan sebuah batas di antaranya adalah sebuah
sakelar, bukan sebuah lereng.
\textbf{18.} $\mathrm{d}\ln\lambda = \tfrac{1}{2}(0.02 - 0.06) = -0.02$
tiap derajat: $\lambda$ turun \qty{2}{\%} tiap derajat, \qty{10}{\%} bagi
\qty{5}{\celsius}, menjadi \qty{99}{\micro m}; batas \emph{hunchback}-nya
bergerak dari $278$ ke \qty{250}{\micro m}, yakni tiga inti ke depan.
Lalat berkembang antara \qty{18}{} dan \qty{29}{\celsius} dengan perbandingan yang normal,
sehingga ada yang mengimbanginya --- sebuah sumber atau pembacaan yang
bergeser bersama suhunya ke arah sebaliknya.
\textbf{19.} $2^{8} = 256$ sandi; empat belas yang dipakai. Dengan keunggulan
belakang, hanya gen paling belakang yang menyala yang berarti, sehingga paling banyak sembilan
jati diri yang berbeda (delapan gennya, atau tidak satu pun): sandinya sebuah peringkat, bukan
sebuah kata biner.
\textbf{20.} Segmen toraks ketiganya menjadi segmen kedua: \emph{Ubx}
menekan acara sayap yang dijalankan sebuah segmen toraks secara bawaan,
dan menyingkirkan penekannya melepaskan keadaan nenek moyangnya --- yakni empat sayap,
seperti pada nenek moyang lalatnya yang bersayap empat.
\textbf{21.} Daerah leher angsanya berjalan sampai ruas tulang belakang 22: yakni sebuah
leher panjang beranggota banyak ruas, sedangkan mencitnya tujuh. Gennya sama;
batas depan pengungkapannya bergerak ke belakang --- yakni sebuah perubahan
pengatur pada tempat sebuah \omterm{def:b3:developmental-genetics:hox}{gen Hox} dinyalakan.
\textbf{22.} Tiap tepinya kini sebuah sumber; Shh-nya paling tinggi pada kedua tepinya
(jari 4), turun ke arah tengahnya (3, lalu 2), lalu tidak pernah mencapai aras
rendah yang menentukan jari 1, sehingga tengahnya memuat dua jari 2 yang saling
membelakangi: 4-3-2-2-3-4.
\textbf{23.} \emph{Pax6} adalah sebuah pengatur induk yang lestari yang dapat
memicu acara mata apa pun yang diusung inangnya: lalatnya membangun sebuah mata
lalat, sehingga gen mencitnya tidak mengusung keterangan mata mencit, melainkan hanya
perintah ``bangunlah sebuah mata di sini.'' Ia menunjukkan \omterm{def:b3:developmental-genetics:evodevo}{homologi dalam} jaringan
gennya, bukan homologi organnya: mata majemuk dan mata kamera
berevolusi terpisah dari sebuah nenek moyang yang punya sepetak fotoreseptor
\emph{Pax6}.
\textbf{24.} Lima jari melintasi \qty{2}{mm} adalah sebuah panjang gelombang
\qty{0.4}{mm}; enam memerlukan \qty{0.33}{mm}, yakni sebuah pemendekan seperenam.
Karena $\lambda \propto \sqrt{D_{\text{inh}}/k_{\text{inh}}}$, naikkanlah
peluruhan penghambatnya sebesar \qty{44}{\%} atau turunkanlah pembaurannya sebesar
\qty{31}{\%}; pada lempeng tangannya, mengurangi dosis \omterm{def:b3:developmental-genetics:hox}{gen Hox} belakangnya
melakukannya, dan mencit mutannya punya enam, tujuh, atau lebih jari yang
tipis.
\textbf{25.} $\lambda \approx \qty{110}{\micro m}$; sebuah penggandaan
sumbernya menggeser tiap batasnya sejauh \qty{76}{\micro m}; derau \qty{10}{\%}
pada kepekatannya adalah \qty{11}{\micro m}, yakni sekitar satu inti, galat
kedudukannya.
\end{solution}
